\section{Macchine di Turing}

\begin{itemize}
    \item $M = (Q, \Sigma, \delta, s)$: Macchina di Turing deterministica.
    \item $N = (Q, \Sigma, \nu, s)$: Macchina di Turing non deterministica.
    \item $k \in \mathbb{N}$: Numero di nastri.
\end{itemize}

\hrulefill

\begin{itemize}
    \item $Q$: Insieme finito di stati.
    \item $s \in Q$: Stato iniziale.
    \item $"yes" \in Q$: Stato di accettazione.
    \item $"no" \in Q$: Stato di rigetto.
    \item $h \in Q$: Stato di terminazione.
\end{itemize}

\hrulefill

\begin{itemize}
    \item $\Sigma$: Alfabeto (insieme finito di simboli).
    \item $\underline{\sigma} \in \Sigma$: Simbolo sotto la testina di un nastro.
    \item $\triangleright$: Start.
    \item $\sqcup$: Blank.
    \item $\Sigma_\sqcup = \Sigma \cup \{\sqcup\}$
    \item $\Gamma = \Sigma \cup \{\triangleright, \sqcup\}$
    %\item $L \subseteq \Sigma^*$: Linguaggio.
\end{itemize}

\hrulefill

\begin{itemize}
    \item $\delta: Q \times \Sigma^k \rightarrow Q \times \Sigma^k \times D^k$ funzione di transizione: prende in input i simboli sotto le testine e lo stato corrente, restituisce simboli che sovrascrivono quelli letti, uno stato in cui transitare al passo successivo e gli spostamenti che ogni testina deve effettuare effettuare alla fine.
    \item $D = \{\leftarrow, \rightarrow, -\}$ spostamenti della testina sul nastro: puó muoversi di una cella a sinistra o a destra, oppure restare ferma.
    \item $\delta(q, \underline{\triangleright}) = (q', \triangleright, \rightarrow)$ (con $q, q' \in Q$) se la testina si trova su $\triangleright$, deve riscriverlo e spostarsi obbligatoriamente a destra.
\end{itemize}

\hrulefill

\begin{itemize}
    \item $x \in \Sigma^*$: Input.
    \item $n = |x| \in \mathbb{N}$
    \item $f(n)$: Tempo utilizzato dalla computazione.
    \item $g(n)$: Spazio di memoria utilizzato dalla computazione.
\end{itemize}

\hrulefill

% \begin{itemize}
%     \item $U$: Macchina di Turing Universale
%     \item $\underline{M}$: Codifica di $M$
%     \item $X = (\underline{M}, x)$: Input di $U$
%     \item $N = |X|$
% \end{itemize}

% \hrulefill

\subsection{Definizione di configurazione}
Siano, $\forall i \in \{1, \dots, k\}$:
\begin{itemize}
    \item $q \in Q$: Stato corrente.
    \item $w_i = "\triangleright w_{i,2} w_{i,3} \dots \underline{w_{i,{m_i}}}" \in \Gamma^*$: Stringa sul nastro $i$ prima della testina, incluso il simbolo sotto la testina.
    (Nota: all'inizio della computazione, $w_i = \underline{\triangleright}$).
    \item $u_i = "u_{i,1} u_{i,2} \dots u_{i,{p_i}}" \in {\Sigma_\cup}^*$: Stringa sul nastro $i$ dopo la testina.
\end{itemize}

\[
\configurazione{q}{w}{u}
\]

\subsubsection{Massimo numero di configurazioni possibili}

\begin{itemize}
    \item Numero di stati: $|Q|$
    \item Massimo numero di stringhe possibili: $|\Sigma_\cup|^{g(n) \cdot k}$ \\
    ($|\Sigma_\cup|^{g(n)}$ stringhe per ogni nastro)
    \item Massimo numero di posizioni delle testine possibili: $g(n)^k$ \\
    ($g(n)$ posizioni per ogni testina)
\end{itemize}

{\[|Q| \cdot |{\Sigma_\cup|}^{g(n) \cdot k} \cdot g(n)^k\]}


\hrulefill

\subsection{Definizione di computazione}
\subsubsection{Computazione di un singolo passo}
\[
\configurazione{q}{w}{u}
\xrightarrow{M}
\configurazione{q'}{w'}{u'}
\]

Se e solo se: \\

Se $\delta(q, [\underline{w_{i{m_i}}}]_{i=1 \dots k}) = (q', \overline{\sigma}, \overline{d})$, dove:
\begin{itemize}
    \item $q' \in Q$: Stato successivo della macchina.
    \item $\overline{\sigma} = (\sigma_i, \dots, \sigma_k) \in {\Sigma_\cup}^k$: Vettore di k caratteri, ognuno da scrivere su una delle $k$ celle sotto una testina.
    \item $\overline{d} = (d_i, \dots, d_k) \in D^k$: Vettore di k direzioni che ogni testina deve seguire.
\end{itemize}

Allora, $\forall i \in \{1, \dots, k\} :$
{\small
\begin{enumerate}
    \item $(d_i = -) \implies ((w'_i = "\triangleright w_{i,2} w_{i,3} \dots \underline{\sigma_i}") \land (u'_i = u_i))$
    \item $(d_i = \leftarrow) \implies ((w'_i = "\triangleright w_{i,2} w_{i,3} \dots \underline{w_{i,m-1}}") \land (u'_i = "\sigma_i u_{i,1} u_{i,2} \dots u_{i,p_i}"))$
    \item $((d_i = \rightarrow) \land (u_i \neq \epsilon)) \implies ((w'_i = "\triangleright w_{i,2} w_{i,3} \dots \sigma_i \underline{u_{i,1}}") \land (u'_i = "u_{i,2} \dots u_{i,p_i}")) $
    \item $((d_i = \rightarrow) \land (u_i = \epsilon)) \implies ((w'_i = "\triangleright w_{i,2} w_{i,3} \dots \sigma_i \underline{\sqcup}") \land (u'_i = \epsilon)) $
\end{enumerate}
}

\hrulefill

\subsubsection{Computazione di $t$ passi}

\[
\configurazione{q_0}{w_0}{u_0}
\xrightarrow{M^t}
\configurazione{q_t}{w_t}{u_t}
\] \\

Se $\forall j \in \{0, \dots, t-1\} :$
\[
\configurazione{q_j}{w_j}{u_j}
\xrightarrow{M}
\configurazione{q_{j+1}}{w_{j+1}}{u_{j+1}}
\]

\hrulefill

\subsubsection{Computazione}

\[
\configurazione{q}{w}{u}
\xrightarrow{M^*}
\configurazione{q'}{w'}{u'}
\] \\

Se $\exists t \in \mathbb{N} :$
\[
\configurazione{q}{w}{u}
\xrightarrow{M^t}
\configurazione{q'}{w'}{u'}
\]

\hrulefill

\subsection{Output}

    $M(x) \in \Sigma^* \cup \{"yes", "no",\nearrow$\}
    \begin{itemize}
        \item $M(x) = "yes"$ sse $\exists w, u :$
        \[
        \configurazionestart{x} \xrightarrow{M^*} \configurazione{"yes"}{w}{u}
        \]
        \item $M(x) = "no"$ sse $\exists w, u :$
        \[
        \configurazionestart{x} \xrightarrow{M^*} \configurazione{"no"}{w}{u}
        \]
        \item $M(x) = y \in \Sigma^*$ sse $\exists w, u :$
        \[
        \configurazionestart{x} \xrightarrow{M^*} \configurazione{h}{w}{u} \land 
        y = "w_k u_k"
        \]
        \item $M(x) = \nearrow$ se la macchina diverge (la computazione non termina). Cioé se non vale nessuno dei casi precedenti.
	\end{itemize}
    
\hrulefill

\subsection{Linguaggi e funzioni}

\subsubsection{Linguaggi ricorsivi}
 $L \subseteq \Sigma^*$ é \textbf{ricorsivo} sse \[\exists M :  M(x) =
	\begin{cases}
		"yes" \qquad se \qquad x \in L \\
		"no" \qquad altrimenti
    \end{cases}
   \] \\
   $M$ \textbf{decide} $L$. \\
   
   \subsubsection{Linguaggi ricorsivamente enumerabili}
    $L \subseteq \Sigma^*$ é \textbf{ricorsivamente enumerabile} sse \[\exists M :  M(x) =
   \begin{cases}
   	"yes" \qquad se \qquad x \in L \\
   	\nearrow \qquad altrimenti
   \end{cases}
   \] \\
   $M$ \textbf{accetta} $L$. \\
   
   \subsubsection{Proposizione 2.1}
   \boxed{L \quad \textbf{ricorsivo} \implies L \quad \textbf{ricorsivamente enumerabile}} \\
   \textbf{Prova:} \\
   \[\exists M :  M(x) =
   \begin{cases}
   	"yes" \qquad se \qquad x \in L \\
   	"no" \qquad altrimenti
   \end{cases}
   \] \\
   Costruiamo $M' = (Q', \Sigma, \delta', s)$ aggiungendo un nuovo stato $q'$ che, nella nuova funzione di transizione $\delta'$, prende il posto dello stato terminale $"no"$. La funzione di transizione é uguale a quella di M in tutti gli altri casi. Il nuovo stato $q'$ sposta la testina a destra all'infinito. Formalmente:
   \begin{itemize}
   		\item $Q' = Q \cup \{q'\}$ \quad dove \quad $q' \notin Q$
   		\item $\delta(q, \sigma) = ("no",  \tau, d) \implies \delta'(q, \sigma) = (q', \tau, d)  \qquad \forall q \in Q;  \forall \sigma, \tau \in \Gamma;  \forall d \in D$
   		\item $\delta(q, \sigma) = (p,  \tau, d) \implies \delta'(q, \sigma) = (p, \tau, d)  \quad \forall q, p \in Q \setminus \{"no"\};  \forall \sigma, \tau \in \Gamma;  \forall d \in D$
   		\item $\delta'(q', \sigma) = (q', \sigma, \rightarrow) \qquad \forall \sigma \in \Sigma \cup \{\sqcup\}$
   \end{itemize}
      \[M'(x) =
   \begin{cases}
   	"yes" \qquad se \qquad x \in L \\
   	\nearrow \qquad altrimenti
   \end{cases}
   \] \\
   $QED$
   
   \subsubsection{Funzioni ricorsive}
   La funzione totale $f : \Sigma^* \rightarrow {\Sigma_\sqcup}^*$ é \textbf{ricorsiva} sse \[\exists M :  \forall x \in \Sigma^* : M(x) = f(x) \] \\
   $M$ \textbf{calcola} $f$. \\
   
      \subsubsection{Funzioni ricorsivamente enumerabili}
   La funzione parziale $f : \Sigma^* \rightarrow {\Sigma_\sqcup}^*$ é \textbf{ricorsivamente enumerabile} sse \[\exists M :  \forall x \in \Sigma^* : \begin{cases}
   	M(x) = y \qquad se \quad f(x) = y \\
   	M(x) = \nearrow \qquad altrimenti
   \end{cases}
    \] \\
    Ricordando che le funzioni parziali non sono definite su tutto il dominio, e dunque esistono stringhe $x \in \Sigma^*$ per cui $f$ non é definita.
    
 \subsection{Complessitá di tempo}
 $M(x)$ richiede tempo $t$ sse:
 \[
	 \configurazionestart{x} \xrightarrow{M^t} \configurazione{H}{w}{u}
 \]
 Dove $H \in \{h, "yes", "no"\}. $\\
 
\textbf{$M$ opera in tempo $f(n)$ sse $\forall x \in \Sigma^*$ $M(x)$ richiede tempo $t <= f(n)$}

\subsection{Complessitá di spazio}
\subsubsection{Macchine di Turing con input e output}
Per definire la complessitá di spazio di memoria di un problema, occorre prima introdurre il concetto di macchine di Turing con input e output.
\begin{itemize}
	\item $k >= 2$
	\item Il primo nastro non viene mai modificato (read only).
	\item L'ultimo nastro viene scandito solo verso destra (write only).
	\item Il cursore del primo nastro non supera mai il primo $\sqcup$ successivo all'input $x$.
\end{itemize}
\subsubsection{Teorema: potenza delle MdT con input e output}
Per ogni macchina di Turing con $k$ nastri $M$ che opera in tempo $f(n)$, esiste una macchina di Turing con input e output e $k+2$ nastri $M'$ che ha lo stesso output e opera in tempo $O(f(n))$. \\
\textbf{Prova:} M' copia l'input sul secondo nastro, simula M sui k nastri centrali, infine, se M ha un output, lo copia sull'ultimo nastro.